% 12-appendix.tex — Reproducibility and Evidence-Level Coding
% Included via % 12-appendix.tex — Reproducibility and Evidence-Level Coding
% Included via % 12-appendix.tex — Reproducibility and Evidence-Level Coding
% Included via % 12-appendix.tex — Reproducibility and Evidence-Level Coding
% Included via \input{sections/12-appendix} in main.tex (after 11-conclusion)

\appendix

\section{Reproducibility and Evidence-Level Coding}
\label{sec:appendix-evidence}

\subsection{Evidence-Level Legend}
\label{sec:appendix-legend}

Claims about commercial and closed-source systems throughout
\S\ref{sec:systems} are annotated with evidence-quality markers that
indicate the epistemic basis for each architectural assertion.  The four
levels, printed as small superscripts in the text, are:

\begin{description}[leftmargin=2em, labelwidth=1.6em, nosep, font=\normalfont]
  \item[\texttt{[D]}] \textbf{Official documentation.}
    Claim is directly supported by the vendor's published API
    documentation, technical report, model card, or developer guide.
    This is the strongest evidence class for a closed-source system.

  \item[\texttt{[S]}] \textbf{Source code inspected.}
    Claim is verified by direct inspection of a publicly available
    repository (GitHub, PyPI source distribution, etc.).  Applicable
    primarily to open-source and developer-tool systems.

  \item[\texttt{[V]}] \textbf{Vendor-reported.}
    Claim originates from a blog post, press release, product
    changelog, or marketing material published by the vendor.
    Vendor-reported claims are reproducible via the cited URL but are
    not independently verified and may reflect marketing rather than
    precise engineering description.

  \item[\texttt{[I]}] \textbf{Inferred from behavior.}
    Claim is derived from observable system behavior during testing,
    community reverse-engineering reports, or logical inference from
    observable inputs and outputs.  No official confirmation exists.
    This is the weakest evidence class; such claims should be treated
    as informed hypotheses rather than established facts.
\end{description}

\subsection{Evidence Assignment Protocol}
\label{sec:appendix-protocol}

Evidence levels were assigned according to the following protocol.
Where multiple sources support a claim, the highest-quality level is
reported: \texttt{[S]} $>$ \texttt{[D]} $>$ \texttt{[V]} $>$ \texttt{[I]}.

\begin{enumerate}[leftmargin=*, nosep]
  \item \textbf{Open-source systems} (Mem0, Zep/Graphiti, Letta,
    HippoRAG, SYNAPSE, LightMem, and all systems with public
    repositories) receive \texttt{[S]} by default, as architectural
    claims are verifiable against source code.  These systems are not
    annotated inline since the evidence basis is uniform.

  \item \textbf{Commercial consumer platforms} (ChatGPT Memory, Gemini
    Memory, Grok Memory, Meta AI Memory) receive \texttt{[D]} when
    official API or help-center documentation exists, \texttt{[V]} for
    blog and changelog claims, and \texttt{[I]} for behavioral
    observations without official confirmation.

  \item \textbf{Developer tools with partial source availability}
    (Claude API Memory Tool, Claude Code auto-memory) receive
    \texttt{[S]} where the relevant code is inspectable and \texttt{[D]}
    for documented API behavior.

  \item Claims whose evidence basis changed between drafts retain the
    evidence level of the most recent source consulted.  The survey
    cutoff is May 2026.
\end{enumerate}

\subsection{Per-System Coding Record}
\label{sec:appendix-matrix}

Table~\ref{tab:system-matrix} provides a tabular overview of all 150
surveyed systems, recording each system's RAG tier (Naive, Advanced,
Modular, or File-Based) and eight discriminating capability features
(Episodic, Semantic, Graph, Fusion, Consolidation, Forgetting,
Contradiction resolution, Preference modeling).  These eight binary
features were selected because they are (a) verifiable from public
sources and (b) the most frequently cited discriminators in the
literature.  The matrix does \emph{not} encode the full four-axis coding
scheme---working/procedural memory, lifecycle sub-categories, and
preference subtypes (N/E/T/M) are captured in the structured data
accompanying this paper.  The complete per-system coding records in
machine-readable JSON and CSV format are available in the project
repository~\cite{alaya2026}.  Evidence-level annotations for commercial
systems appear inline in \S\ref{sec:production} of the main text; they
are not reproduced in the matrix because the matrix encodes
architectural features, not evidence quality.

\subsection{Benchmark Evaluation Code and Data}
\label{sec:appendix-code}

All baseline evaluation code (LoCoMo, LongMemEval\_S, and MAB
replications), raw per-question results, and judge logs are available in
the project repository~\cite{alaya2026} under \texttt{bench/}.  The
system coding data for the 148 systems with complete matrix records is
provided as:
\begin{itemize}[leftmargin=*, nosep]
  \item \texttt{survey-paper/data/coding-records.csv} --- RAG tier and
    eight binary capability features for all 133 systems.
  \item \texttt{survey-paper/data/coding-records.json} --- same data in
    JSON with schema metadata.
  \item \texttt{survey-paper/data/intercoder-reliability.yaml} ---
    dual-coded sample (28 systems), per-axis $\kappa$ values,
    and adjudication log for all 12 disagreements.
\end{itemize}
See the repository \texttt{README} for exact reproduction steps,
dependency versions, and the \texttt{all-MiniLM-L6-v2} model checkpoint
hash used for na\"ive RAG embeddings.
 in main.tex (after 11-conclusion)

\appendix

\section{Reproducibility and Evidence-Level Coding}
\label{sec:appendix-evidence}

\subsection{Evidence-Level Legend}
\label{sec:appendix-legend}

Claims about commercial and closed-source systems throughout
\S\ref{sec:systems} are annotated with evidence-quality markers that
indicate the epistemic basis for each architectural assertion.  The four
levels, printed as small superscripts in the text, are:

\begin{description}[leftmargin=2em, labelwidth=1.6em, nosep, font=\normalfont]
  \item[\texttt{[D]}] \textbf{Official documentation.}
    Claim is directly supported by the vendor's published API
    documentation, technical report, model card, or developer guide.
    This is the strongest evidence class for a closed-source system.

  \item[\texttt{[S]}] \textbf{Source code inspected.}
    Claim is verified by direct inspection of a publicly available
    repository (GitHub, PyPI source distribution, etc.).  Applicable
    primarily to open-source and developer-tool systems.

  \item[\texttt{[V]}] \textbf{Vendor-reported.}
    Claim originates from a blog post, press release, product
    changelog, or marketing material published by the vendor.
    Vendor-reported claims are reproducible via the cited URL but are
    not independently verified and may reflect marketing rather than
    precise engineering description.

  \item[\texttt{[I]}] \textbf{Inferred from behavior.}
    Claim is derived from observable system behavior during testing,
    community reverse-engineering reports, or logical inference from
    observable inputs and outputs.  No official confirmation exists.
    This is the weakest evidence class; such claims should be treated
    as informed hypotheses rather than established facts.
\end{description}

\subsection{Evidence Assignment Protocol}
\label{sec:appendix-protocol}

Evidence levels were assigned according to the following protocol.
Where multiple sources support a claim, the highest-quality level is
reported: \texttt{[S]} $>$ \texttt{[D]} $>$ \texttt{[V]} $>$ \texttt{[I]}.

\begin{enumerate}[leftmargin=*, nosep]
  \item \textbf{Open-source systems} (Mem0, Zep/Graphiti, Letta,
    HippoRAG, SYNAPSE, LightMem, and all systems with public
    repositories) receive \texttt{[S]} by default, as architectural
    claims are verifiable against source code.  These systems are not
    annotated inline since the evidence basis is uniform.

  \item \textbf{Commercial consumer platforms} (ChatGPT Memory, Gemini
    Memory, Grok Memory, Meta AI Memory) receive \texttt{[D]} when
    official API or help-center documentation exists, \texttt{[V]} for
    blog and changelog claims, and \texttt{[I]} for behavioral
    observations without official confirmation.

  \item \textbf{Developer tools with partial source availability}
    (Claude API Memory Tool, Claude Code auto-memory) receive
    \texttt{[S]} where the relevant code is inspectable and \texttt{[D]}
    for documented API behavior.

  \item Claims whose evidence basis changed between drafts retain the
    evidence level of the most recent source consulted.  The survey
    cutoff is May 2026.
\end{enumerate}

\subsection{Per-System Coding Record}
\label{sec:appendix-matrix}

Table~\ref{tab:system-matrix} provides a tabular overview of all 150
surveyed systems, recording each system's RAG tier (Naive, Advanced,
Modular, or File-Based) and eight discriminating capability features
(Episodic, Semantic, Graph, Fusion, Consolidation, Forgetting,
Contradiction resolution, Preference modeling).  These eight binary
features were selected because they are (a) verifiable from public
sources and (b) the most frequently cited discriminators in the
literature.  The matrix does \emph{not} encode the full four-axis coding
scheme---working/procedural memory, lifecycle sub-categories, and
preference subtypes (N/E/T/M) are captured in the structured data
accompanying this paper.  The complete per-system coding records in
machine-readable JSON and CSV format are available in the project
repository~\cite{alaya2026}.  Evidence-level annotations for commercial
systems appear inline in \S\ref{sec:production} of the main text; they
are not reproduced in the matrix because the matrix encodes
architectural features, not evidence quality.

\subsection{Benchmark Evaluation Code and Data}
\label{sec:appendix-code}

All baseline evaluation code (LoCoMo, LongMemEval\_S, and MAB
replications), raw per-question results, and judge logs are available in
the project repository~\cite{alaya2026} under \texttt{bench/}.  The
system coding data for the 148 systems with complete matrix records is
provided as:
\begin{itemize}[leftmargin=*, nosep]
  \item \texttt{survey-paper/data/coding-records.csv} --- RAG tier and
    eight binary capability features for all 133 systems.
  \item \texttt{survey-paper/data/coding-records.json} --- same data in
    JSON with schema metadata.
  \item \texttt{survey-paper/data/intercoder-reliability.yaml} ---
    dual-coded sample (28 systems), per-axis $\kappa$ values,
    and adjudication log for all 12 disagreements.
\end{itemize}
See the repository \texttt{README} for exact reproduction steps,
dependency versions, and the \texttt{all-MiniLM-L6-v2} model checkpoint
hash used for na\"ive RAG embeddings.
 in main.tex (after 11-conclusion)

\appendix

\section{Reproducibility and Evidence-Level Coding}
\label{sec:appendix-evidence}

\subsection{Evidence-Level Legend}
\label{sec:appendix-legend}

Claims about commercial and closed-source systems throughout
\S\ref{sec:systems} are annotated with evidence-quality markers that
indicate the epistemic basis for each architectural assertion.  The four
levels, printed as small superscripts in the text, are:

\begin{description}[leftmargin=2em, labelwidth=1.6em, nosep, font=\normalfont]
  \item[\texttt{[D]}] \textbf{Official documentation.}
    Claim is directly supported by the vendor's published API
    documentation, technical report, model card, or developer guide.
    This is the strongest evidence class for a closed-source system.

  \item[\texttt{[S]}] \textbf{Source code inspected.}
    Claim is verified by direct inspection of a publicly available
    repository (GitHub, PyPI source distribution, etc.).  Applicable
    primarily to open-source and developer-tool systems.

  \item[\texttt{[V]}] \textbf{Vendor-reported.}
    Claim originates from a blog post, press release, product
    changelog, or marketing material published by the vendor.
    Vendor-reported claims are reproducible via the cited URL but are
    not independently verified and may reflect marketing rather than
    precise engineering description.

  \item[\texttt{[I]}] \textbf{Inferred from behavior.}
    Claim is derived from observable system behavior during testing,
    community reverse-engineering reports, or logical inference from
    observable inputs and outputs.  No official confirmation exists.
    This is the weakest evidence class; such claims should be treated
    as informed hypotheses rather than established facts.
\end{description}

\subsection{Evidence Assignment Protocol}
\label{sec:appendix-protocol}

Evidence levels were assigned according to the following protocol.
Where multiple sources support a claim, the highest-quality level is
reported: \texttt{[S]} $>$ \texttt{[D]} $>$ \texttt{[V]} $>$ \texttt{[I]}.

\begin{enumerate}[leftmargin=*, nosep]
  \item \textbf{Open-source systems} (Mem0, Zep/Graphiti, Letta,
    HippoRAG, SYNAPSE, LightMem, and all systems with public
    repositories) receive \texttt{[S]} by default, as architectural
    claims are verifiable against source code.  These systems are not
    annotated inline since the evidence basis is uniform.

  \item \textbf{Commercial consumer platforms} (ChatGPT Memory, Gemini
    Memory, Grok Memory, Meta AI Memory) receive \texttt{[D]} when
    official API or help-center documentation exists, \texttt{[V]} for
    blog and changelog claims, and \texttt{[I]} for behavioral
    observations without official confirmation.

  \item \textbf{Developer tools with partial source availability}
    (Claude API Memory Tool, Claude Code auto-memory) receive
    \texttt{[S]} where the relevant code is inspectable and \texttt{[D]}
    for documented API behavior.

  \item Claims whose evidence basis changed between drafts retain the
    evidence level of the most recent source consulted.  The survey
    cutoff is May 2026.
\end{enumerate}

\subsection{Per-System Coding Record}
\label{sec:appendix-matrix}

Table~\ref{tab:system-matrix} provides a tabular overview of all 150
surveyed systems, recording each system's RAG tier (Naive, Advanced,
Modular, or File-Based) and eight discriminating capability features
(Episodic, Semantic, Graph, Fusion, Consolidation, Forgetting,
Contradiction resolution, Preference modeling).  These eight binary
features were selected because they are (a) verifiable from public
sources and (b) the most frequently cited discriminators in the
literature.  The matrix does \emph{not} encode the full four-axis coding
scheme---working/procedural memory, lifecycle sub-categories, and
preference subtypes (N/E/T/M) are captured in the structured data
accompanying this paper.  The complete per-system coding records in
machine-readable JSON and CSV format are available in the project
repository~\cite{alaya2026}.  Evidence-level annotations for commercial
systems appear inline in \S\ref{sec:production} of the main text; they
are not reproduced in the matrix because the matrix encodes
architectural features, not evidence quality.

\subsection{Benchmark Evaluation Code and Data}
\label{sec:appendix-code}

All baseline evaluation code (LoCoMo, LongMemEval\_S, and MAB
replications), raw per-question results, and judge logs are available in
the project repository~\cite{alaya2026} under \texttt{bench/}.  The
system coding data for the 148 systems with complete matrix records is
provided as:
\begin{itemize}[leftmargin=*, nosep]
  \item \texttt{survey-paper/data/coding-records.csv} --- RAG tier and
    eight binary capability features for all 133 systems.
  \item \texttt{survey-paper/data/coding-records.json} --- same data in
    JSON with schema metadata.
  \item \texttt{survey-paper/data/intercoder-reliability.yaml} ---
    dual-coded sample (28 systems), per-axis $\kappa$ values,
    and adjudication log for all 12 disagreements.
\end{itemize}
See the repository \texttt{README} for exact reproduction steps,
dependency versions, and the \texttt{all-MiniLM-L6-v2} model checkpoint
hash used for na\"ive RAG embeddings.
 in main.tex (after 11-conclusion)

\appendix

\section{Reproducibility and Evidence-Level Coding}
\label{sec:appendix-evidence}

\subsection{Evidence-Level Legend}
\label{sec:appendix-legend}

Claims about commercial and closed-source systems throughout
\S\ref{sec:systems} are annotated with evidence-quality markers that
indicate the epistemic basis for each architectural assertion.  The four
levels, printed as small superscripts in the text, are:

\begin{description}[leftmargin=2em, labelwidth=1.6em, nosep, font=\normalfont]
  \item[\texttt{[D]}] \textbf{Official documentation.}
    Claim is directly supported by the vendor's published API
    documentation, technical report, model card, or developer guide.
    This is the strongest evidence class for a closed-source system.

  \item[\texttt{[S]}] \textbf{Source code inspected.}
    Claim is verified by direct inspection of a publicly available
    repository (GitHub, PyPI source distribution, etc.).  Applicable
    primarily to open-source and developer-tool systems.

  \item[\texttt{[V]}] \textbf{Vendor-reported.}
    Claim originates from a blog post, press release, product
    changelog, or marketing material published by the vendor.
    Vendor-reported claims are reproducible via the cited URL but are
    not independently verified and may reflect marketing rather than
    precise engineering description.

  \item[\texttt{[I]}] \textbf{Inferred from behavior.}
    Claim is derived from observable system behavior during testing,
    community reverse-engineering reports, or logical inference from
    observable inputs and outputs.  No official confirmation exists.
    This is the weakest evidence class; such claims should be treated
    as informed hypotheses rather than established facts.
\end{description}

\subsection{Evidence Assignment Protocol}
\label{sec:appendix-protocol}

Evidence levels were assigned according to the following protocol.
Where multiple sources support a claim, the highest-quality level is
reported: \texttt{[S]} $>$ \texttt{[D]} $>$ \texttt{[V]} $>$ \texttt{[I]}.

\begin{enumerate}[leftmargin=*, nosep]
  \item \textbf{Open-source systems} (Mem0, Zep/Graphiti, Letta,
    HippoRAG, SYNAPSE, LightMem, and all systems with public
    repositories) receive \texttt{[S]} by default, as architectural
    claims are verifiable against source code.  These systems are not
    annotated inline since the evidence basis is uniform.

  \item \textbf{Commercial consumer platforms} (ChatGPT Memory, Gemini
    Memory, Grok Memory, Meta AI Memory) receive \texttt{[D]} when
    official API or help-center documentation exists, \texttt{[V]} for
    blog and changelog claims, and \texttt{[I]} for behavioral
    observations without official confirmation.

  \item \textbf{Developer tools with partial source availability}
    (Claude API Memory Tool, Claude Code auto-memory) receive
    \texttt{[S]} where the relevant code is inspectable and \texttt{[D]}
    for documented API behavior.

  \item Claims whose evidence basis changed between drafts retain the
    evidence level of the most recent source consulted.  The survey
    cutoff is May 2026.
\end{enumerate}

\subsection{Per-System Coding Record}
\label{sec:appendix-matrix}

Table~\ref{tab:system-matrix} provides a tabular overview of all 150
surveyed systems, recording each system's RAG tier (Naive, Advanced,
Modular, or File-Based) and eight discriminating capability features
(Episodic, Semantic, Graph, Fusion, Consolidation, Forgetting,
Contradiction resolution, Preference modeling).  These eight binary
features were selected because they are (a) verifiable from public
sources and (b) the most frequently cited discriminators in the
literature.  The matrix does \emph{not} encode the full four-axis coding
scheme---working/procedural memory, lifecycle sub-categories, and
preference subtypes (N/E/T/M) are captured in the structured data
accompanying this paper.  The complete per-system coding records in
machine-readable JSON and CSV format are available in the project
repository~\cite{alaya2026}.  Evidence-level annotations for commercial
systems appear inline in \S\ref{sec:production} of the main text; they
are not reproduced in the matrix because the matrix encodes
architectural features, not evidence quality.

\subsection{Benchmark Evaluation Code and Data}
\label{sec:appendix-code}

All baseline evaluation code (LoCoMo, LongMemEval\_S, and MAB
replications), raw per-question results, and judge logs are available in
the project repository~\cite{alaya2026} under \texttt{bench/}.  The
system coding data for the 148 systems with complete matrix records is
provided as:
\begin{itemize}[leftmargin=*, nosep]
  \item \texttt{survey-paper/data/coding-records.csv} --- RAG tier and
    eight binary capability features for all 133 systems.
  \item \texttt{survey-paper/data/coding-records.json} --- same data in
    JSON with schema metadata.
  \item \texttt{survey-paper/data/intercoder-reliability.yaml} ---
    dual-coded sample (28 systems), per-axis $\kappa$ values,
    and adjudication log for all 12 disagreements.
\end{itemize}
See the repository \texttt{README} for exact reproduction steps,
dependency versions, and the \texttt{all-MiniLM-L6-v2} model checkpoint
hash used for na\"ive RAG embeddings.
